% experiments.tex
% Experiments and Results section for VICTOR paper.
% Intended to be \input{} from the main paper .tex file, after methodology.tex.

% \gap{...} marks a remaining open item that still needs a decision or a
% citation before submission. Search for "\gap{" to find every open item.
% Defined with \providecommand so this file is safe to \input from a main
% paper .tex that may also define \gap in its own preamble.
\providecommand{\gap}[1]{\textbf{[GAP: #1]}}

\section{Experiments and Results}
\label{sec:experiments}

We evaluate \method on \NumTasks{} meta-training tasks (\NumEpisodes{}
episodes, \NumTrainHours{} hours of teleoperated bimanual data) and a
held-out suite of \NumTaskPairs{} in-domain/out-of-domain task pairs
(\NumEvalTasks{} tasks total) on the YOR bimanual manipulator~\cite{yor}, in
five stages: (1) value-function and retrieval quality, (2) the main policy
comparison, (3) a demonstration-count ablation, (4) a retrieval--policy
mismatch and context ablation, and (5) a context-sensitivity analysis.

\subsection{Value Function and Retrieval Quality}
\label{sec:vfe-retrieval-quality}

\begin{table}[t]
    \centering
    \caption{\textbf{Value-function and retrieval quality (placeholder).}
    Value correlation is measured against human progress annotations on the
    \NumEvalTasks{} evaluation tasks; retrieval MAE uses each method's value
    predictions in place of the IC-VFE to drive retrieval
    (Section~\ref{sec:vfe-retrieval-quality}).}
    \label{tab:vfe-retrieval-results}
    \footnotesize
    \begin{tabular}{lrrrr}
        \toprule
        & \multicolumn{2}{c}{Value correlation} & \multicolumn{2}{c}{Retrieval MAE} \\
        Method & Kendall $\tau_b$ & Pearson $r$ & Position & Orientation \\
        \midrule
        Robometer (zero-shot) & X.XX & X.XX & X.XX & X.XX \\
        TOPReward             & X.XX & X.XX & X.XX & X.XX \\
        GVL                   & X.XX & X.XX & X.XX & X.XX \\
        RECAP (fine-tuned)    & X.XX & X.XX & X.XX & X.XX \\
        Robometer (fine-tuned) & X.XX & X.XX & X.XX & X.XX \\
        IC-VFE (ours)         & \textbf{X.XX} & \textbf{X.XX} & \textbf{X.XX} & \textbf{X.XX} \\
        \midrule
        RoboDopamine\textsuperscript{$\dagger$} & X.XX & X.XX & X.XX & X.XX \\
        \bottomrule
    \end{tabular}
    \\[2pt]
    \raggedright\footnotesize\textsuperscript{$\dagger$}The annotation model itself, fine-tuned as a value estimator; a sanity check on how well its own training target can be fit, not a compared baseline.
\end{table}

We compare zero-shot value function estimators (Robometer~\cite{robometer},
TOPReward~\cite{topreward}, GVL~\cite{gvl}) against estimators fine-tuned on
our full meta-training dataset with dense RoboDopamine annotations
(RECAP~\cite{pi06}, Robometer, and our IC-VFE). Each method is scored by (a)
Kendall's $\tau_b$ and Pearson correlation against human-annotated progress
values, and (b) our retrieval-quality metric: using a method's value
predictions in place of the IC-VFE to retrieve a context chunk $c=(o',a')$
via $f_{\text{retrieve}}$ (Section~\ref{sec:context-retrieval}), then
computing the MAE between $a'$ and the ground-truth action paired with the
query observation, separately for position and orientation. Both metrics for
all methods appear in Table~\ref{tab:vfe-retrieval-results}.

\gap{[Analysis]}

\subsection{Main Evaluation}
\label{sec:main-eval}

\begin{figure}[t]
    \centering
    \includegraphics[width=0.95\linewidth]{figures/main_sr_barchart.pdf}
    \caption{\textbf{Total success rate by method (placeholder).} Bars show
    in-domain and out-of-domain success rate (\%) for each compared method,
    over \NumRolloutsPerTask{} rollouts per task across all \NumEvalTasks{}
    tasks. \gap{placeholder figure -- replace with real results.}}
    \label{fig:main-sr}
\end{figure}

We compare $\pi_{0.5}$ trained for 40k and 50k steps (no retrieval),
RICL-$\pi_{0.5}$~\cite{ricl}, and three \method variants: value-based
retrieval, vision-similarity retrieval, and vision+value retrieval (our main
method). All methods are trained on the same meta-training data and
evaluated over \NumRolloutsPerTask{} rollouts per task on all \NumEvalTasks{}
tasks. We report total and partial (subtask-level) success rate, split by
in-domain and out-of-domain task, in Figure~\ref{fig:main-sr}.

\gap{[Analysis]}

\subsection{Demonstration-Count Ablation}
\label{sec:demo-count-ablation}

\begin{figure}[t]
    \centering
    \includegraphics[width=0.95\linewidth]{figures/demo_count_sr.pdf}
    \caption{\textbf{Success rate vs.\ number of context demonstrations
    (placeholder).} $x$-axis $\in \{1,3,5,7,10\}$ demonstrations, $y$-axis
    success rate (\%), pooled over 6 tasks. \gap{placeholder figure --
    replace with real results.}}
    \label{fig:demo-count}
\end{figure}

Taking the best policy from Section~\ref{sec:main-eval}, we vary the number
of demonstrations available in the retrieval buffer over $\{1, 3, 5, 7, 10\}$,
randomly selected from the available pool, on 6 long-horizon tasks (3
in-domain and their 3 out-of-domain counterparts), with 10 rollouts per
configuration ($6 \times 10 \times 5 = 300$ rollouts total).

\gap{[Analysis]}

\subsection{Retrieval--Policy Mismatch and Context Ablation}
\label{sec:mismatch-ablation}

\begin{table}[t]
    \centering
    \caption{\textbf{Train/test retrieval mismatch (placeholder).} Success
    rate (\%) when the retrieval strategy at test time differs from the one
    used during meta-training.}
    \label{tab:retrieval-mismatch}
    \footnotesize
    \begin{tabular}{lrr}
        \toprule
        Train $\backslash$ Test & Value retrieval & Vision retrieval \\
        \midrule
        Value retrieval  & XX.X & XX.X \\
        Vision retrieval & XX.X & XX.X \\
        \bottomrule
    \end{tabular}
    \quad
    \begin{tabular}{lrr}
        \toprule
        Method & Default & Swapped \\
        \midrule
        $\pi_{0.5}$ (+ context) & XX.X & XX.X \\
        \method (no context)    & XX.X & XX.X \\
        \bottomrule
    \end{tabular}
\end{table}

We measure how much retrieval choices actually matter at test time in two
ways. First, we cross the training-time and test-time retrieval strategy
(value vs.\ vision) for our two single-signal \method variants in a
$2\times2$ design. Second, we swap context presence: $\pi_{0.5}$ (trained
without context) is given a retrieved context chunk at test time, and
\method (trained with context) is deployed with no context at test time.
Both results are reported as success rate in Table~\ref{tab:retrieval-mismatch}.

\gap{[Analysis]}

\subsection{Context Sensitivity}
\label{sec:context-sensitivity}

\begin{table}[t]
    \centering
    \caption{\textbf{Context sensitivity (placeholder).} MAE between the
    action predicted from the true retrieved context and the action
    predicted from random context chunks; higher indicates a policy relies
    more on its retrieved context.}
    \label{tab:context-sensitivity}
    \footnotesize
    \begin{tabular}{lrr}
        \toprule
        Method & Position MAE & Orientation MAE \\
        \midrule
        RICL-$\pi_{0.5}$        & X.XX & X.XX \\
        Ours (value retrieval)  & X.XX & X.XX \\
        Ours (vision retrieval) & X.XX & X.XX \\
        Ours (vision + value)   & X.XX & X.XX \\
        \bottomrule
    \end{tabular}
\end{table}

For each policy, we take an observation, retrieve its context chunk via that
policy's own retrieval strategy, and record the predicted action $a_c$; we
then replace the context with randomly sampled chunks from the demonstration
pool and record the resulting action $a_{\text{rand}}$. We report the MAE
between $a_c$ and $a_{\text{rand}}$, separately for position and
orientation, in Table~\ref{tab:context-sensitivity}.

\gap{[Analysis]}

